\documentclass{article}
\usepackage[left=2cm, right=2cm, top=2cm, bottom=2cm]{geometry}
\usepackage{graphicx}
\usepackage{amsmath}
\usepackage{amssymb}
\usepackage{amsthm}
\usepackage{fancyhdr}
\usepackage{verbatim}
\usepackage{listings}
\usepackage{xcolor}
\usepackage{pdfpages}
\usepackage{cancel}
\usepackage{physics}
\usepackage{esint}

\lstset{
    language=C++,
    basicstyle=\ttfamily\footnotesize,
    keywordstyle=\color{blue},
    commentstyle=\color{green},
    stringstyle=\color{red},
    numbers=left,
    numberstyle=\tiny\color{gray},
    frame=single,
    breaklines=true,
    captionpos=b,
}


\title{MTH 553: Midterm \#2 Recap}
\author{Cliff Sun}

\newtheorem{theorem}{Theorem}[section]
\newtheorem{lemma}[theorem]{Lemma}
\newtheorem{definition}[theorem]{Definition}
\newtheorem{conjecture}[theorem]{Conjecture}
\newtheorem{proposition}[theorem]{Proposition}
\newtheorem{corollary}[theorem]{Corollary}
\newtheorem{one minute paper}[theorem]{One Minute Paper}

\pagestyle{fancy}
\lhead{\textbf{\thepage}\ \ \nouppercase{\rightmark}}
\chead{MTH 553: Midterm \#2 Recap}
\rhead{Cliff Sun}

\begin{document}

\maketitle

\section*{Key points}
\begin{enumerate}
    \item Green's Formulas
    \item Some questions may be stated in $\mathbb{R}^2$ or $\mathbb{R}^3$
    \item Lower order terms in heat equation (change of variables and reduce to heat equation)
    \item Lower order terms in wave equation (change of variables and reduce to 2d wave equation)
\end{enumerate}

\textbf{Green's Formulas:} Let $\Omega$ be some domain. Then (Green 1) states that 
\begin{equation}
    \int_{\Omega}u \triangle v dx + \int_{\Omega}\nabla u \cdot \nabla v dx = \int_{\partial \Omega}u\frac{\partial v}{\partial \nu}dS
\end{equation}

Here, $\nu$ is the normal vector to $\partial \Omega$. Green 2 states that 
\begin{equation}
    \int_{\Omega}\left( u\triangle v - v \triangle u \right)dx = \int_{\partial\Omega}\left( u\frac{\partial v}{\partial \nu} - v\frac{\partial u}{\partial \nu} \right)dx
\end{equation}

Here, $v$ has compact support and is infinitely differentiable. Lower order terms:

\begin{equation}
    u_t = k \triangle u + \vec{b} \cdot \nabla u+ cu + f(x,t)
\end{equation}
Here, $k$ is the diffusion rate, $\vec{b}$ is the transport term (direction of propagation or moving coordinate frame), $c$ is the heat generation term, and $f(x,t)$ is an external source. To figure out what the 
terms mean, just set all the other terms to be zero and reason out the equation. 

For lower order terms in the heat equation, use change of variables. For the wave equation with the lower order terms, change of variables and solve with 2d wave equation.

Continuous dependence for the heat equation. Using maximum principle, $|u_1 - u_2| \leq \epsilon$ initially for all $x$ and $|u_1 - u_2| \leq \epsilon$ on $\partial \Omega$. Then the max principle on $w = u_1 - u_2$ and $-w$ implies that $|w| \leq \epsilon$ on $\Omega_T$. 
For enegry method, then $|u_1 - u_2| \leq \epsilon$ initially and $u_1 - u_2 \equiv 0$ on $\partial \Omega$. So then $w \equiv 0$ on $\partial \Omega$. Then you obtain $E'(t) \leq 0$. 

Mininmum principle is just the maximum principle but applied to $-u$. That is, $u$ achieves its minimum value somewhere on the parabolic boundary. You can also create a strong minimum principle.  
\end{document}